\section{Base}


\section{Joints}


\section{Arms}


\section{End-Effector}
The end-effector is the moving plate connected to the end of the arms. It's primary function is to hold the modules which allow the robot to interact with the world around it.

\subsection{Structure}
To allow for high acceleration at lower power cost, the structure's inertia must be kept at a minimum.
Lower inertia allows for higher acceleration at lower forces, since inertia is directly linked to an object's mass.
Thus the main goal for the structural design is minimizing weight, while keeping the end-effector stable enough to hold heavier loads without breaking. \\
A rounded triangular model was chosen, since it's three sides allow seamless integration with the three robot arms, while minimizing material around the edges.


\section{Modules}